\documentclass[a4paper,9pt]{extarticle}
\usepackage[UTF8]{ctex}
\usepackage[top=16mm,bottom=14mm,left=12mm,right=12mm,columnsep=7mm]{geometry}
\usepackage{fancyhdr,titlesec,xcolor,amsmath,amssymb,listings,graphicx}
\usepackage{needspace,tabularx,array,booktabs}
\usepackage[colorlinks=true,linkcolor=accent,urlcolor=accent,bookmarksnumbered=true]{hyperref}
\definecolor{accent}{HTML}{24485D}
\definecolor{codebg}{HTML}{F5F7F8}
\definecolor{muted}{HTML}{61717B}
\IfFontExistsTF{Times New Roman}{\setmainfont{Times New Roman}}{}
\IfFontExistsTF{Arial}{\setsansfont{Arial}}{}
\IfFontExistsTF{Consolas}{\setmonofont{Consolas}[Scale=0.9]}{\setmonofont{Latin Modern Mono}}
\lstset{language=C++,basicstyle=\fontsize{7.6}{9.2}\selectfont\ttfamily,
  keywordstyle=\bfseries\color{accent},commentstyle=\color{muted},
  stringstyle=\color{black!80},numbers=left,numberstyle=\fontsize{5.5}{7}\selectfont\color{muted},
  numbersep=5pt,backgroundcolor=\color{codebg},
  showspaces=false,showstringspaces=false,showtabs=false,tabsize=2,
  columns=flexible,keepspaces=true,breaklines=true,breakatwhitespace=false,
  breakindent=1em,
  frame=l,framerule=0.4pt,rulecolor=\color{accent!35},framesep=3pt,
  aboveskip=5pt,belowskip=5pt,xleftmargin=10pt,xrightmargin=3pt}
\pagestyle{fancy}\fancyhf{}
\fancyhead[L]{\fontsize{8}{10}\selectfont\sffamily\color{muted}ACM / ICPC}
\fancyhead[R]{\fontsize{8}{10}\selectfont\sffamily\color{muted}\leftmark}
\fancyfoot[C]{\fontsize{8}{10}\selectfont\sffamily\thepage}
\renewcommand{\sectionmark}[1]{\markboth{#1}{}}
\renewcommand{\headrulewidth}{0.3pt}\setlength{\headheight}{13pt}
\titleformat{\section}{\large\bfseries\sffamily\color{accent}}{\thesection}{0.6em}{}[\titlerule]
\titleformat{\subsection}{\normalsize\bfseries\sffamily\color{accent}}{\thesubsection}{0.5em}{}
\titleformat{\subsubsection}{\small\bfseries}{\thesubsubsection}{0.5em}{}
\titlespacing*{\section}{0pt}{10pt}{5pt}
\titlespacing*{\subsection}{0pt}{8pt}{3pt}
\titlespacing*{\subsubsection}{0pt}{5pt}{2pt}
\setcounter{tocdepth}{2}\setcounter{secnumdepth}{2}
\setlength{\parindent}{0pt}\setlength{\parskip}{2.5pt}
\setlength{\emergencystretch}{2em}
\raggedbottom
\newsavebox{\mathbox}
\newcommand{\fitmath}[1]{\par\smallskip\begingroup
  \sbox{\mathbox}{$\displaystyle #1$}
  \noindent\makebox[\linewidth][c]{\ifdim\wd\mathbox>\linewidth
    \resizebox{\linewidth}{!}{\usebox{\mathbox}}\else\usebox{\mathbox}\fi}
  \par\endgroup\smallskip}
\hypersetup{pdftitle={ACM/ICPC 算法模板},pdfauthor={Rejected}}
\begin{document}

\thispagestyle{empty}
\vspace*{12mm}
{\Huge\bfseries\sffamily\color{accent}ACM / ICPC}\par
\vspace{3mm}{\Large\sffamily 算法模板手册}\par
\vspace{3mm}{\color{muted}Rejected\enspace |\enspace 2026-10-11}\par
\vspace{5mm}\hrule\vspace{4mm}
1 个章节 · 11 个条目 · 11 个代码块\par\medskip


\Needspace{5\baselineskip}
\subsection*{使用约定}

本手册各代码块独立选用；同名结构体、多个 main 和不同版本不要一起粘贴。片段按题意补齐输入输出与全局数组。

\noindent\textbullet\enspace 默认使用 GNU C++17；\texttt{\_\_int128}、\texttt{bits/\allowbreak{}stdc+\allowbreak{}+\allowbreak{}.\allowbreak{}h} 和 builtin 属于 GCC 扩展，个别片段另需 C+\allowbreak{}+\allowbreak{}20/\allowbreak{}23。\par
\noindent\textbullet\enspace 常用别名：\texttt{using ll = long long; using db = double; using i128 = \_\_int128\_t; using pii = pair<\allowbreak{}int,\allowbreak{}int>\allowbreak{};}。\par
\noindent\textbullet\enspace 下标和区间以各条目说明为准；多测清空边表、计数器、懒标记与答案，递归深树注意栈空间。\par
\noindent\textbullet\enspace 先估算时间与内存；乘法在计算前转宽类型，模逆先确认存在，哈希和浮点结果注意碰撞与误差。\par
\noindent\textbullet\enspace 原有题目解法、样例、重复版本全部保留；标为“片段/骨架”的代码须补齐依赖后使用。\par

\Needspace{5\baselineskip}
\subsection*{如何添加算法}

在对应一级章节下添加二级算法标题，可在代码前写“用法 / 注意 / 复杂度”。一个条目支持多个代码块、正文、公式和表格。

\begin{lstlisting}[language={}]
## 算法名称

**用法**：调用入口、下标和返回值。**注意**：适用条件。

```cpp
// 你的代码
```
\end{lstlisting}

运行 \texttt{python make.py} 生成 PDF；\texttt{python make.py --no-pdf} 仅生成 TeX。需要 XeLaTeX（MiKTeX / TeX Live），Python 无需第三方包。

版式参考：\href{https://github.com/kth-competitive-programming/kactl}{KACTL}、\href{https://github.com/jaehyunp/stanfordacm}{Stanford ACM Notebook}。

\par\smallskip

\clearpage\twocolumn
\tableofcontents
\clearpage
\Needspace{5\baselineskip}
\section{不要忘记}

\Needspace{5\baselineskip}
\subsection{土制cph(仅linux)}

\textbf{用法}：保存为 Python 脚本，传入题号和源文件名；读取 samples-题号中的 .\allowbreak{}in/\allowbreak{}.\allowbreak{}ans。\textbf{注意}：依赖 Linux timeout 与 g+\allowbreak{}+；文件名不要含空格，先确认编译成功。

\begin{lstlisting}[language={Python}]
from os import system as e, listdir as l
from sys import argv
_, q, f = argv  # q: 题号, f: 文件名

if "py" in f:
	r = f"python3 {f}"
else:
	e(f"g++ -std=c++23 -O2 -Wall {f}.cpp -o {f}")
	r = "./" + f

#
d = "samples-" + q.capitalize()
for i in l(d):
	if not "in" in i: continue
	print(i)
	p = d + "/" + i[:-2]
	if e(f"timeout 2 {r}<{p}in>{p}out"): print("TLE or RE")
	else:
		with open(f"{p}out") as o, open(f"{p}ans") as a:
			print("AC" if o.read().split()==a.read().split() else "WA")
\end{lstlisting}

\Needspace{5\baselineskip}
\subsection{土制cph(全平台)}

\textbf{用法}：同上，Python 分支传带 .py 的文件名，C++ 分支传不含扩展名的文件名。\textbf{注意}：需可用的 Python/\allowbreak{}g+\allowbreak{}+，非零退出码按 RE 处理；浮点比较须同时检查 token 数量。

\begin{lstlisting}[language={Python}]
from os import listdir as l, name as osname
from os.path import join as jp
from sys import argv as a, executable
import subprocess as s
_,q,f = a
if f.endswith(".py"):
    cmd = [executable, f]
else:
    exe = f + (".exe" if osname == "nt" else "")
    s.run(["g++", "-std=gnu++20", "-O2", "-Wall", f + ".cpp", "-o", exe], check=True)
    cmd = [jp(".", exe)]

d='samples-'+q.capitalize()
g,r,p,n='\033[32m','\033[31m','\033[35m','\033[0m'
for i in l(d):
    if not i.endswith('.in'): continue
    t=jp(d,i[:-2])
    print(i,'测评结果 ',end='',flush=1)
    try:
        k = s.run(cmd, timeout=2, stdin=open(f"{t}in"), 
                 stdout=open(f"{t}out", 'w'), stderr=s.PIPE, text=1)
        if k.returncode:
            print(f'{r}RE{n}')
            if k.stderr: print(k.stderr)
            continue
        c = lambda o,a : o.read().split() == a.read().split()
        # c = lambda o,a : all(abs(float(x)-float(y))/max(1,abs(float(y)))<=1e-4 
        #               for x,y in zip(o.read().split(),a.read().split()))
        print(f'{g}AC{n}' if c(open(f"{t}out"),open(f"{t}ans")) else f'{r}WA{n}')
        if k.stderr:print(f'{r}{k.stderr}{n}')
    except s.TimeoutExpired as k:
        print(f'{p}TLE or RE{n}')
        if k.stderr: print(f'{r}{k.stderr}{n}')
\end{lstlisting}

\Needspace{5\baselineskip}
\subsection{随机数}

\textbf{用法}：Rand(x,y) 返回闭区间 [x,y] 的均匀随机整数，要求 x<=y；shuffle 打乱序列。\textbf{注意}：对拍复现时固定种子。

\begin{lstlisting}[language={C++}]
mt19937_64 rng(chrono::steady_clock::now().time_since_epoch().count());
int Rand(int x, int y) {
	return uniform_int_distribution<int>(x, y)(rng);
}
shuffle(a.begin(), a.end(), rng);
\end{lstlisting}

\Needspace{5\baselineskip}
\subsection{对拍}

\textbf{用法}：依次运行生成器、暴力程序和待测程序，比较 A.out 与 A.ans。\textbf{注意}：先确认三个命令都成功；Windows 用 fc，Linux 用 diff，保留首个失败输入。

\begin{lstlisting}[language={C++}]
// checker.cpp
while (1) {
	system("A_gen > _data.in");  // linux 下换成 ./xxx
	system("A_bf < _data.in > A.ans");
	system("A < _data.in > A.out");
	if (system("fc A.out A.ans")) {  // linux 下换成 diff
		cout << "WA\n"; return 0;
	} else {
		cout << "AC\n";
	}
}
\end{lstlisting}

\Needspace{5\baselineskip}
\subsection{解除python大数限制}

\textbf{用法}：Python 3.11+ 在转换超长十进制整数前设置；此处仍限制为 100005 位。\textbf{注意}：set\_\allowbreak{}int\_\allowbreak{}max\_\allowbreak{}str\_\allowbreak{}digits(\allowbreak{}0) 才是完全解除；旧版本需判断该函数是否存在。

\begin{lstlisting}[language={Python}]
import sys
sys.set_int_max_str_digits(100005)
\end{lstlisting}

\Needspace{5\baselineskip}
\subsection{常用函数}

\textbf{用法}：read 读取 int，qpow(\allowbreak{}a,\allowbreak{}b) 求模幂，qpow(a) 求逆元。\textbf{注意}：快读假定输入完整；模运算传入 [0,MOD) 的数，逆元要求 MOD 为质数且 a 非零；乘积超 ll 时用 i128。

\begin{lstlisting}[language={C++}]
// 快读
inline int read()
{
	int x=0,f=1;char ch=getchar();
	while (ch<'0'||ch>'9'){if (ch=='-') f=-1;ch=getchar();}
	while (ch>='0'&&ch<='9'){x=x*10+ch-48;ch=getchar();}
	return x*f;
}

// 取模
const int MOD=998244353;
ll add(ll x, ll y) { return (x+y)%MOD; }
ll del(ll x, ll y) { return add(x, MOD-y); }
ll mul(ll x, ll y) { return (x*y)%MOD; }
ll qpow(ll a, ll b = MOD - 2) {
	ll res = 1;
	for (; b; b >>= 1) {
		if (b & 1) { res = mul(res, a); }
		a = mul(a, a);
	}
	return res;
}

// 计算一个整数的二进制表示中有多少个 1
__builtin_popcountll(i);
std::popcount((unsigned long long)i);  // c++20

// 计算一个整数的二进制表示中最高位的 1
__lg(i);  // i > 0；按实参类型支持 int / long long
\end{lstlisting}

\Needspace{5\baselineskip}
\subsection{高精度}

\textbf{用法}：先定义 \_\_int128 的流运算符，再以 cin/\allowbreak{}cout 读写；fixed/\allowbreak{}setprecision 只控制浮点显示。\textbf{注意}：\_\_int128 是固定 128 位整数，不是任意精度；本版不处理最小负数取反。

\begin{lstlisting}[language={C++}]

cout << fixed << setprecision(15);//输出精度

istream& operator>> (istream& is, __int128 &x)
{
	string s; is>>s;
	bool f=false;
	__int128 ans = 0;
	for (auto c: s) {
		if (c == '-') { f = true; }
		if (isdigit(c)) { ans = ans * 10 + (c-'0'); }
	}
	x = f? -ans: ans;
	return is;
}
ostream& operator<< (ostream& os, __int128 &x)
{
	string ans;
	function<void(__int128)> write = [&](__int128 x){
		if (x < 0) { ans += '-'; x=-x; }
		if (x > 9) { write(x/10); }
		ans += '0'+x%10;
	};
	write(x); return os << ans;
}
\end{lstlisting}
\Needspace{5\baselineskip}
\subsection{多进程评测}

\textbf{用法}：此处仅为接收进程编号并读写数据的示例。\textbf{注意}：它不会自动启动多个进程；并发调度与每个进程的独立输入输出需另写。

\begin{lstlisting}[language={C++}]
#include <cstdlib>
#include <iostream>
using namespace std;
int main(int argc, char *argv[]) {
    int id = argc > 1 ? atoi(argv[1]) : 0;
    int v;
    cin >> v;
    cout << v << endl;
    cout.flush();
    return 0;
}
\end{lstlisting}
\Needspace{5\baselineskip}
\subsection{迭代器辅助函数}

\textbf{注意}：不要递增 end()、递减 begin() 或解引用 end()；set 的迭代器不支持减法，distance 是线性时间。
\begin{lstlisting}[language={C++}]

函数	作用	是否修改原迭代器	set 上复杂度
std::next(it)	返回 it 后一个位置的迭代器	否	O(1)
std::next(it, n)	返回 it 后 n 步的迭代器	否	O(n)
std::prev(it)	返回 it 前一个位置的迭代器	否	O(1)
std::prev(it, n)	返回 it 前 n 步的迭代器	否	O(n)
std::advance(it, n)	把 it 本身移动 n 步	是	O(n)
std::distance(a, b)	计算 a 到 b 的步数	否	O(n)
++it / --it	前进一步 / 后退一步	是	O(1)
注意：
std::next、std::prev 返回的是新迭代器，原来的 it 不变。

std::advance 会直接改 it，没有返回值。

对 set 来说，std::distance 是 O(n)，不要在高频查询中使用。元素个数与 key 差值不同；要 O(log n) 查排名，请使用树状数组/平衡树。
\end{lstlisting}
\Needspace{5\baselineskip}
\subsection{set相关}

\textbf{用法}：优先用 s.\allowbreak{}lower\_\allowbreak{}bound/\allowbreak{}upper\_\allowbreak{}bound 进行对数查找。\textbf{注意}：比较器须严格弱序，不能用 <\allowbreak{}=；multiset.\allowbreak{}erase(\allowbreak{}x) 会删掉全部 x，只删一个时用 find 后 erase(it)。
\begin{lstlisting}[language={C++}]
// 默认：从小到大（std::less<int>）
set<int> s1;
// 从大到小（std::greater<int>）
set<int, greater<int>> s2;
// 自定义比较器（仿函数）
struct GreaterCmp {
    bool operator()(int a, int b) const {
        return a > b;   // 从大到小
    }
};
set<int, GreaterCmp> s3;
// 用 lambda 自定义比较器（注意：lambda 类型没有默认构造函数）
auto cmp = [](int a, int b) { return a > b; };
set<int, decltype(cmp)> s4(cmp);
// insert(value) 返回 pair<iterator, bool>
//   - first: 指向插入元素或已存在元素
//   - second: 是否插入成功
// insert(hint, value) 返回 iterator
// hint 正确时均摊 O(1)，否则 O(log n)
s1.insert(s1.end(), 20);
// 插入范围
vector<int> v = {5, 6, 7};
s1.insert(v.begin(), v.end());
// emplace 原地构造
s1.emplace(30);
// ------------------------------------------------------------
// 3. 删除
// ------------------------------------------------------------
// 按 key 删除，返回删除个数（set 中为 0 或 1）
size_t cnt = s1.erase(10);
auto it2 = s1.find(20);
if (it2 != s1.end()) {
    it2 = s1.erase(it2);   // it2 现在指向下一个元素
}

// 删除范围 [first, last)
// s1.erase(s1.begin(), s1.end());
// 清空
// s1.clear();
// ------------------------------------------------------------
// 4. 查找
// ------------------------------------------------------------
// find(key) 找到返回迭代器，否则返回 end()
auto it3 = s1.find(5);
if (it3 != s1.end()) {
    // 找到了
}
// count(key) set 中返回 0 或 1
if (s1.count(5)) {
    // 存在
}

// lower_bound(key) / upper_bound(key)
// 默认 less 时：
//   lower_bound: 第一个 >= key
//   upper_bound: 第一个 > key
// 使用 greater 时语义相反（降序）：
//   lower_bound: 第一个 <= key
//   upper_bound: 第一个 < key
auto lb = s1.lower_bound(5);
auto ub = s1.upper_bound(5);

// equal_range 返回 pair(lower_bound, upper_bound)
auto [lo, hi] = s1.equal_range(5);
// ------------------------------------------------------------
// 5. 迭代器与容量
// ------------------------------------------------------------
// s1.begin(), s1.end()
// s1.cbegin(), s1.cend()
// s1.rbegin(), s1.rend()
// s1.size(), s1.empty()
// ------------------------------------------------------------

\end{lstlisting}

\Needspace{5\baselineskip}
\subsection{vector相关}

\textbf{注意}：min/\allowbreak{}max\_\allowbreak{}element 要求非空；unique 只去相邻重复，须先排序；reserve 不改变 size，resize 会；扩容后旧迭代器、引用可能失效。

\begin{lstlisting}[language={C++}]
// 多维 vector，需要 c++17，c++14 只能老老实实写 vector<vector<int>>
vector a(n+1, vector (n+1, vector<int>(n+1)));

// 最值
int mx = *max_element(vec.begin(), vec.end(), cmp);
int mn = *min_element(vec.begin(), vec.end(), cmp);

// 去重与离散化
for (int i = 1; i <= n; ++i) { vec.push_back(a[i]); }
sort(vec.begin(), vec.end());
vec.erase(unique(vec.begin(), vec.end()), vec.end());
for (int i = 1; i <= n; ++i) {
	idx[i]=lower_bound(vec.begin(), vec.end(), a[i])-vec.begin()+1;  // 下标从1开始
}

// 前缀和
vector <ll> a(n), s(n);
partial_sum(a.begin(), a.end(), s.begin());
ll sum = accumulate(a.begin(), a.end(), 0LL);

// 填充 0~n-1，令 a[i]=i
iota(a.begin(), a.end(), 0);
\end{lstlisting}

\end{document}